% =============================================================================
% Algorithm walkthrough section for the HR visual-search paper.
%
% Required packages:
%   amsmath, amssymb, mathtools
%   algorithm, algpseudocode
%
% The file is a section, not a standalone document. Paste it into the paper or
% use % =============================================================================
% Algorithm walkthrough section for the HR visual-search paper.
%
% Required packages:
%   amsmath, amssymb, mathtools
%   algorithm, algpseudocode
%
% The file is a section, not a standalone document. Paste it into the paper or
% use % =============================================================================
% Algorithm walkthrough section for the HR visual-search paper.
%
% Required packages:
%   amsmath, amssymb, mathtools
%   algorithm, algpseudocode
%
% The file is a section, not a standalone document. Paste it into the paper or
% use % =============================================================================
% Algorithm walkthrough section for the HR visual-search paper.
%
% Required packages:
%   amsmath, amssymb, mathtools
%   algorithm, algpseudocode
%
% The file is a section, not a standalone document. Paste it into the paper or
% use \input{algorithm_walkthrough_section} after loading the packages above.
% =============================================================================

\providecommand{\indicator}[1]{\mathbf{1}\!\left[#1\right]}
\providecommand{\set}[1]{\mathcal{#1}}
\providecommand{\argmax}{\operatorname*{arg\,max}}
\providecommand{\argmin}{\operatorname*{arg\,min}}
\providecommand{\proc}[1]{\text{\textsc{#1}}}

\section{Algorithm Walkthrough}
\label{sec:algorithm-walkthrough}

\subsection{High-level idea}

The algorithm is a bounded, best-first tree search over image regions. It does
not build the complete tree in advance. Instead, it starts from the full image
and repeatedly selects one frontier region, evaluates it, and, if the region is
promising and has not reached a stopping rule, splits it into smaller child
regions.

Each iteration therefore has four steps:
\begin{enumerate}
  \item \textbf{Select.} Choose one frontier node using a cost-aware UCT rule.
  \item \textbf{Evaluate.} Run the detector on the selected crop, verify the
        candidate boxes with the VLM, and compute a region reward.
  \item \textbf{Back up.} Add the reward to every node on the selected path.
  \item \textbf{Expand.} Split the evaluated node into child regions, or mark it
        exhausted when no useful split exists.
\end{enumerate}

The intuition is simple. A region that receives a high reward raises the
average reward of its whole branch, so that branch is more likely to be selected
again. At the same time, unvisited children are always explored first, and a UCT
exploration bonus favours branches with few visits. A cost penalty discourages
branches whose remaining work is expensive. The search is therefore an
exploitation--exploration trade-off, not a hard threshold on the VLM score.

The search outputs a selected region \(R\), a set of candidate boxes
\(\set{B}\), and the full tree and allocation trace. In the complete system, the
final free-form answer is generated by the VLM from the selected region:
\begin{equation}
  a = \proc{Answer}(I_R, q),
\end{equation}
where \(I_R\) is the image crop induced by \(R\). The answer is a downstream
product; the tree reward is only a localization heuristic.

% -----------------------------------------------------------------------------
\subsection{Search tree and notation}

Let
\begin{equation}
  \set{T} = (\set{V}, \set{E})
\end{equation}
be the search tree. The root \(v_0\) is the full image region
\begin{equation}
  B_{v_0} = (0, 0, W, H).
\end{equation}
Each node \(v \in \set{V}\) represents a region \(B_v\), and stores
\begin{align}
  n_v &\in \mathbb{N}_0
    && \text{completed evaluations backed up through } v, \\
  R_v &\in \mathbb{R}_{\ge 0}
    && \text{sum of raw rewards backed up through } v, \\
  Q_v &=
    \begin{cases}
      R_v / n_v, & n_v > 0, \\
      0, & n_v = 0,
    \end{cases}
    && \text{empirical mean reward}, \\
  \set{C}(v) &\subseteq \set{V}
    && \text{active children of } v.
\end{align}
A node is \emph{exhausted} when it has no unevaluated descendant within the
depth limit. The frontier is the set of active nodes at which a selected path
can end.

The tree is generated lazily. Initially, only the root exists; it is expanded
once to produce its first children. Every other child is created only after its
parent has been completely evaluated and split. This is important: the algorithm
does not enumerate all possible regions first and inspect them afterwards.

% -----------------------------------------------------------------------------
\subsection{Region evaluation}

We describe the candidate-first path, which is the default. The legacy path
scores the whole region before detection. Let \(v\) be the frontier node
selected at iteration \(t\). The detector is run on the crop
\(I_v = I[B_v]\). It returns a set of candidate detections
\begin{equation}
  \set{D}_v
  = \left\{ (b_{v,i}, s^{\mathrm{det}}_{v,i}) \right\}_{i=1}^{m_v},
\end{equation}
where \(b_{v,i}=(x_{1,i},y_{1,i},x_{2,i},y_{2,i})\) is a box in the crop's
coordinate system and \(s^{\mathrm{det}}_{v,i}\in[0,1]\) is the detector
confidence. Low-confidence boxes are removed, and confidence-ordered non-maximum
suppression keeps a set of distinct candidates
\begin{equation}
  \set{D}^{+}_v
  = \left\{ (b_{v,i}, s^{\mathrm{det}}_{v,i}) \right\}_{i=1}^{m^{+}_v}.
\end{equation}
If \(\set{D}^{+}_v=\varnothing\), the VLM is not called:
\begin{equation}
  r_v = 0.
  \label{eq:empty-region-reward}
\end{equation}
This gate is important for efficiency, because most background regions contain
no candidate at all.

Otherwise, the VLM verifies the candidates, by default against the complete
question \(q\). For candidate \(i\), let
\begin{equation}
  y_{v,i} \in \{\text{Yes}, \text{No}, \text{Uncertain}\},
  \qquad
  \hat{s}^{\mathrm{vlm}}_{v,i} \in [0,1]
\end{equation}
be the verifier answer and score. A candidate is \emph{verified} when
\begin{equation}
  z_{v,i}
  = \indicator{y_{v,i}=\text{Yes}}
    \indicator{\hat{s}^{\mathrm{vlm}}_{v,i} \ge \tau},
  \label{eq:verified}
\end{equation}
where \(\tau\) is the acceptance threshold. Define the checked set and the
verified set as
\begin{equation}
  \set{V}_v = \{1,\ldots,m^{+}_v\},
  \qquad
  \set{V}^{+}_v = \{ i \in \set{V}_v : z_{v,i}=1 \}.
\end{equation}

Candidate ranking uses a fused score
\begin{equation}
  \tilde{s}_{v,i}
  = \frac{
      w_{\mathrm{vlm}} \hat{s}^{\mathrm{vlm}}_{v,i}
      + w_{\mathrm{det}} s^{\mathrm{det}}_{v,i}
    }{
      w_{\mathrm{vlm}} + w_{\mathrm{det}}
    },
  \label{eq:candidate-fusion}
\end{equation}
with \(\hat{s}^{\mathrm{vlm}}_{v,i}=0\) when \(y_{v,i}=\text{No}\). The tree
reward is instead a region-level score, defined next.

% -----------------------------------------------------------------------------
\subsection{Region reward}

The region reward combines the strongest verifier judgment, detector confidence,
the number of verified boxes, and their spatial diversity:
\begin{align}
  q^{\star}_v
  &= \max_{i \in \set{V}_v}
     \left\{
       \hat{s}^{\mathrm{vlm}}_{v,i}
       : y_{v,i} \ne \text{No}
     \right\}, \\
  \bar{s}_v
  &= \frac{1}{|\set{V}^{+}_v|}
     \sum_{i \in \set{V}^{+}_v} s^{\mathrm{det}}_{v,i}, \\
  \kappa_v
  &= \frac{\min(|\set{V}^{+}_v|, K)}{K}, \\
  \bar{\iota}_v
  &= \frac{1}{\binom{|\set{V}^{+}_v|}{2}}
     \sum_{\{i,j\}\subseteq \set{V}^{+}_v}
     \operatorname{IoU}(b_{v,i}, b_{v,j}), \\
  \delta_v
  &= 1 - \bar{\iota}_v,
\end{align}
where \(K\) is the count cap and \(q^{\star}_v=0\) if no candidate has a
non-\text{No} answer. For \(|\set{V}^{+}_v|<2\), set \(\bar{\iota}_v=0\) and
\(\delta_v=1\). Then
\begin{equation}
  r_v =
  \begin{cases}
    \displaystyle
    \frac{
      w_q q^{\star}_v
      + w_s \bar{s}_v
      + w_\kappa \kappa_v
      + w_\delta \delta_v
    }{
      w_q + w_s + w_\kappa + w_\delta
    }, & \set{V}^{+}_v \ne \varnothing, \\[8pt]
    0, & \set{V}^{+}_v = \varnothing.
  \end{cases}
  \label{eq:region-reward}
\end{equation}
The weights are nonnegative and their sum is positive. With the default
implementation they sum to one, but the normalization in
\eqref{eq:region-reward} keeps the reward in \([0,1]\) even if the weights are
changed.

The reward is a heuristic, not a calibrated probability. In particular,
\(q^{\star}_v\) lets one strong candidate dominate, and the diversity term is
maximally favourable when only one box is verified.

% -----------------------------------------------------------------------------
\subsection{Branch selection}

Let \(v\) be an expanded node and let \(\set{C}(v)\) be its active children. If
some child has never been visited, the selector restricts the pool to those
unvisited children. This guarantees that a new region is inspected before the
search commits further to an already promising branch.

If all active children have been visited, the selector computes the cost-aware
UCT index. For each child \(u\in\set{C}(v)\),
\begin{align}
  Q_u &= \frac{R_u}{n_u}, \\
  U_u &= c \sqrt{\frac{\ln(n_v+1)}{n_u}}, \\
  C_u &= \text{remaining-cost forecast below } u, \\
  \mathrm{index}(u)
  &= Q_u + U_u - \lambda \frac{C_u}{T_{\mathrm{ref}}},
\end{align}
and chooses
\begin{equation}
  u^{\star}
  = \argmax_{u \in \set{C}(v)} \mathrm{index}(u).
  \label{eq:selection}
\end{equation}
Ties are broken by the lowest node ID. The three terms have distinct roles:
\begin{itemize}
  \item \(Q_u\) exploits branches that returned high rewards;
  \item \(U_u\) explores branches with few visits;
  \item the last term penalizes branches whose remaining work is expensive.
\end{itemize}
The exploration coefficient is \(c\), the cost weight is \(\lambda\), and
\(T_{\mathrm{ref}}\) normalizes the cost in seconds. Setting \(\lambda=0\)
recovers plain UCT; removing \(U_u\) gives greedy selection.

The remaining-cost forecast is recursive. For a frontier node \(u\),
\begin{equation}
  C_u =
  \begin{cases}
    \max\left(0, \hat{c}(u)\right), & \set{C}(u)=\varnothing, \\[4pt]
    \displaystyle
    \sum_{w\in\set{C}(u)}
      \frac{n_w + 1/|\set{C}(u)|}
           {\sum_{w'\in\set{C}(u)} n_{w'} + 1}
      C_w, & \set{C}(u)\ne\varnothing,
  \end{cases}
  \label{eq:remaining-cost}
\end{equation}
where \(\hat{c}(u)\) estimates the uncached detector, verifier, and split work
at the frontier. The pseudo-visit term \(1/|\set{C}(u)|\) prevents a division by
zero and makes an unvisited set behave like a simple average.

A selected path is built recursively:
\begin{equation}
  P_t = (v_0, v_1, \ldots, v_t),
  \qquad
  v_{\ell+1} = \text{selected child of } v_{\ell},
\end{equation}
until a frontier leaf \(v_t\) is reached. The actual evaluation is performed on
\(v_t\), but the reward is backed up through every node in \(P_t\).

% -----------------------------------------------------------------------------
\subsection{Backup and expansion}

Let \(r_{v_t}\) be the region reward returned by the complete evaluation of the
selected leaf. Every node on the selected path is updated:
\begin{equation}
  n_u \leftarrow n_u + 1,
  \qquad
  R_u \leftarrow R_u + r_{v_t},
  \qquad
  \forall u \in P_t.
  \label{eq:backup}
\end{equation}
The same raw reward is used at every level; there is no discount and no cost
subtraction in the backed-up value. Cost only affects future selection.
Equivalently, the new mean reward is
\begin{equation}
  Q_u \leftarrow
  \frac{n_u Q_u + r_{v_t}}{n_u + 1},
  \qquad
  \forall u \in P_t.
\end{equation}
This is why a promising region increases the selection probability of its entire
branch, not only of the leaf itself.

After the backup, the evaluated node is expanded. Semantic splitting uses a
DINOv3 feature map followed by PCA, SLIC, and K-means:
\begin{equation}
  F_v = \Pi_3\!\left(\phi_{\mathrm{DINO}}(I_v)\right),
\end{equation}
where \(\Pi_3\) projects the patch features to three PCA components. For each
requested SLIC count \(n\in\set{N}\), let \(\ell^{(n)}_v\) be the superpixel
labels and let
\begin{equation}
  H(s) = - \sum_{c} p_{s,c} \log_2 p_{s,c}
\end{equation}
be the colour entropy of superpixel \(s\). The selected count is
\begin{equation}
  n^{\star}
  = \argmin_{n\in\set{N}}
    \frac{1}{|\ell^{(n)}_v|}
    \sum_{s\in \ell^{(n)}_v} H(s).
  \label{eq:slic-entropy}
\end{equation}
K-means then groups superpixels using their mean PCA colour and normalized
centroid:
\begin{equation}
  \mu_s = \left[ \bar{F}_{v,s}, \alpha \bar{p}_s \right],
  \qquad
  k = \min(B, |\ell^{(n^{\star})}_v|),
\end{equation}
where \(B\) is the branching factor and \(\alpha\) is the spatial weight.
Children are the clusters intersected with the parent mask, tightened to their
bounding boxes.

If the resulting children are too few, do not provide a useful zoom, or have
poor geometry, the split falls back to four spatial quadrants. If the node is at
the maximum depth, has too few active pixels, or is too small, it is marked
exhausted and is not split again.

% -----------------------------------------------------------------------------
\subsection{Stopping criteria}

The search stops when any of the following holds:
\begin{equation}
  \begin{aligned}
    &\text{exhausted:} && \text{the root has no unevaluated descendants,} \\
    &\text{api-budget:} && n_{\mathrm{api}} \ge B_{\mathrm{api}}, \\
    &\text{local-budget:} && n_{\mathrm{local}} \ge B_{\mathrm{local}}, \\
    &\text{iterations:} && t \ge T_{\max}, \\
    &\text{match:} && \text{a verified proposal exceeds the match threshold,} \\
    &\text{error:} && \text{a non-recoverable execution error occurred.}
  \end{aligned}
  \label{eq:stop}
\end{equation}
A partial evaluation interrupted by a budget is recorded but is not backed up:
\begin{equation}
  \text{incomplete}(t)
  \quad\Longrightarrow\quad
  \Delta n_u = 0,\ \Delta R_u = 0
  \quad \forall u\in P_t.
\end{equation}
This prevents a failed or truncated evaluation from being treated as a
zero-reward observation.

% -----------------------------------------------------------------------------
\subsection{High-level pseudocode}

\begin{algorithm}[t]
\caption{Bounded cost-aware visual search}
\label{alg:visual-search}
\begin{algorithmic}[1]
\Require image \(I\), question \(q\), configuration \(\theta\)
\Ensure selected region \(R\), candidate boxes \(\set{B}\), tree \(\set{T}\),
         trace \(\tau\)
\State \(s \gets \proc{Prepare}(I,q,\theta)\)
\State \(\set{T} \gets \proc{NewTree}(\proc{Root}(I))\)
\State \(\proc{Expand}(\set{T}.\text{root}, s)\)
\While{not \(\proc{Stop}(\set{T},s)\)}
  \State \(P \gets \proc{SelectPath}(\set{T},s)\)
  \Comment{Eq.~\eqref{eq:selection}}
  \State \(v \gets \text{last}(P)\)
  \State \(r_v \gets \proc{Evaluate}(v,s)\)
  \Comment{Eqs.~\eqref{eq:verified}--\eqref{eq:region-reward}}
  \State \(\proc{Backup}(P,r_v)\)
  \Comment{Eq.~\eqref{eq:backup}}
  \State \(\proc{Expand}(v,s)\)
  \Comment{DINOv3 \(\to\) SLIC \(\to\) K-means, or quadrants}
\EndWhile
\State \(R,\set{B} \gets \proc{Rank}(\set{T})\)
\State \Return \(R,\set{B},\set{T},\tau\)
\end{algorithmic}
\end{algorithm}

The pseudocode keeps the model internals inside \textsc{Evaluate}. The main loop
contains only the four operations that define the search: selection, evaluation,
backup, and expansion. This separation makes the search logic independent of
the particular detector or VLM.

% -----------------------------------------------------------------------------
\subsection{Operations and cost}

Let \(N=T_{\max}\) be the iteration limit, \(D\) the maximum depth, and \(L\) the
maximum number of candidates verified per region. The dominant costs are:
\begin{itemize}
  \item one detector pass per evaluated region, at most \(N\);
  \item at most \(N\) VLM calls in whole-region mode, and at most \(N L\) in
        per-candidate mode, plus repair attempts;
  \item one feature encoding, one SLIC sweep over \(|\set{N}|\) candidate counts,
        and one K-means pass per split;
  \item \(O(|\set{V}|)\) tree memory, with \(|\set{V}|=O(NB)\) in the worst case
        for branching factor \(B\) and depth \(D\).
\end{itemize}
Caches remove repeated work: a node's detector output and split are computed at
most once, and identical verifier inputs can reuse a cached assessment. The tree
can revisit a branch many times, but a node's own complete evaluation runs at
most once; later visits descend into its children.

% -----------------------------------------------------------------------------
\subsection{Summary}

The algorithm can be read as a sequence of simple statements:
\begin{enumerate}
  \item A promising region raises \(Q\) along its path, so the branch is more
        likely to be explored again.
  \item Unvisited children are always considered first, and the UCT bonus
        protects under-visited branches.
  \item The cost penalty removes expensive branches from the top of the queue.
  \item A complete evaluation is the only event that backs up a reward; partial
        or invalid evaluations do not create artificial zeros.
  \item The search returns ranked regions and boxes. The final answer is produced
        downstream by the VLM from the selected region.
\end{enumerate}
 after loading the packages above.
% =============================================================================

\providecommand{\indicator}[1]{\mathbf{1}\!\left[#1\right]}
\providecommand{\set}[1]{\mathcal{#1}}
\providecommand{\argmax}{\operatorname*{arg\,max}}
\providecommand{\argmin}{\operatorname*{arg\,min}}
\providecommand{\proc}[1]{\text{\textsc{#1}}}

\section{Algorithm Walkthrough}
\label{sec:algorithm-walkthrough}

\subsection{High-level idea}

The algorithm is a bounded, best-first tree search over image regions. It does
not build the complete tree in advance. Instead, it starts from the full image
and repeatedly selects one frontier region, evaluates it, and, if the region is
promising and has not reached a stopping rule, splits it into smaller child
regions.

Each iteration therefore has four steps:
\begin{enumerate}
  \item \textbf{Select.} Choose one frontier node using a cost-aware UCT rule.
  \item \textbf{Evaluate.} Run the detector on the selected crop, verify the
        candidate boxes with the VLM, and compute a region reward.
  \item \textbf{Back up.} Add the reward to every node on the selected path.
  \item \textbf{Expand.} Split the evaluated node into child regions, or mark it
        exhausted when no useful split exists.
\end{enumerate}

The intuition is simple. A region that receives a high reward raises the
average reward of its whole branch, so that branch is more likely to be selected
again. At the same time, unvisited children are always explored first, and a UCT
exploration bonus favours branches with few visits. A cost penalty discourages
branches whose remaining work is expensive. The search is therefore an
exploitation--exploration trade-off, not a hard threshold on the VLM score.

The search outputs a selected region \(R\), a set of candidate boxes
\(\set{B}\), and the full tree and allocation trace. In the complete system, the
final free-form answer is generated by the VLM from the selected region:
\begin{equation}
  a = \proc{Answer}(I_R, q),
\end{equation}
where \(I_R\) is the image crop induced by \(R\). The answer is a downstream
product; the tree reward is only a localization heuristic.

% -----------------------------------------------------------------------------
\subsection{Search tree and notation}

Let
\begin{equation}
  \set{T} = (\set{V}, \set{E})
\end{equation}
be the search tree. The root \(v_0\) is the full image region
\begin{equation}
  B_{v_0} = (0, 0, W, H).
\end{equation}
Each node \(v \in \set{V}\) represents a region \(B_v\), and stores
\begin{align}
  n_v &\in \mathbb{N}_0
    && \text{completed evaluations backed up through } v, \\
  R_v &\in \mathbb{R}_{\ge 0}
    && \text{sum of raw rewards backed up through } v, \\
  Q_v &=
    \begin{cases}
      R_v / n_v, & n_v > 0, \\
      0, & n_v = 0,
    \end{cases}
    && \text{empirical mean reward}, \\
  \set{C}(v) &\subseteq \set{V}
    && \text{active children of } v.
\end{align}
A node is \emph{exhausted} when it has no unevaluated descendant within the
depth limit. The frontier is the set of active nodes at which a selected path
can end.

The tree is generated lazily. Initially, only the root exists; it is expanded
once to produce its first children. Every other child is created only after its
parent has been completely evaluated and split. This is important: the algorithm
does not enumerate all possible regions first and inspect them afterwards.

% -----------------------------------------------------------------------------
\subsection{Region evaluation}

We describe the candidate-first path, which is the default. The legacy path
scores the whole region before detection. Let \(v\) be the frontier node
selected at iteration \(t\). The detector is run on the crop
\(I_v = I[B_v]\). It returns a set of candidate detections
\begin{equation}
  \set{D}_v
  = \left\{ (b_{v,i}, s^{\mathrm{det}}_{v,i}) \right\}_{i=1}^{m_v},
\end{equation}
where \(b_{v,i}=(x_{1,i},y_{1,i},x_{2,i},y_{2,i})\) is a box in the crop's
coordinate system and \(s^{\mathrm{det}}_{v,i}\in[0,1]\) is the detector
confidence. Low-confidence boxes are removed, and confidence-ordered non-maximum
suppression keeps a set of distinct candidates
\begin{equation}
  \set{D}^{+}_v
  = \left\{ (b_{v,i}, s^{\mathrm{det}}_{v,i}) \right\}_{i=1}^{m^{+}_v}.
\end{equation}
If \(\set{D}^{+}_v=\varnothing\), the VLM is not called:
\begin{equation}
  r_v = 0.
  \label{eq:empty-region-reward}
\end{equation}
This gate is important for efficiency, because most background regions contain
no candidate at all.

Otherwise, the VLM verifies the candidates, by default against the complete
question \(q\). For candidate \(i\), let
\begin{equation}
  y_{v,i} \in \{\text{Yes}, \text{No}, \text{Uncertain}\},
  \qquad
  \hat{s}^{\mathrm{vlm}}_{v,i} \in [0,1]
\end{equation}
be the verifier answer and score. A candidate is \emph{verified} when
\begin{equation}
  z_{v,i}
  = \indicator{y_{v,i}=\text{Yes}}
    \indicator{\hat{s}^{\mathrm{vlm}}_{v,i} \ge \tau},
  \label{eq:verified}
\end{equation}
where \(\tau\) is the acceptance threshold. Define the checked set and the
verified set as
\begin{equation}
  \set{V}_v = \{1,\ldots,m^{+}_v\},
  \qquad
  \set{V}^{+}_v = \{ i \in \set{V}_v : z_{v,i}=1 \}.
\end{equation}

Candidate ranking uses a fused score
\begin{equation}
  \tilde{s}_{v,i}
  = \frac{
      w_{\mathrm{vlm}} \hat{s}^{\mathrm{vlm}}_{v,i}
      + w_{\mathrm{det}} s^{\mathrm{det}}_{v,i}
    }{
      w_{\mathrm{vlm}} + w_{\mathrm{det}}
    },
  \label{eq:candidate-fusion}
\end{equation}
with \(\hat{s}^{\mathrm{vlm}}_{v,i}=0\) when \(y_{v,i}=\text{No}\). The tree
reward is instead a region-level score, defined next.

% -----------------------------------------------------------------------------
\subsection{Region reward}

The region reward combines the strongest verifier judgment, detector confidence,
the number of verified boxes, and their spatial diversity:
\begin{align}
  q^{\star}_v
  &= \max_{i \in \set{V}_v}
     \left\{
       \hat{s}^{\mathrm{vlm}}_{v,i}
       : y_{v,i} \ne \text{No}
     \right\}, \\
  \bar{s}_v
  &= \frac{1}{|\set{V}^{+}_v|}
     \sum_{i \in \set{V}^{+}_v} s^{\mathrm{det}}_{v,i}, \\
  \kappa_v
  &= \frac{\min(|\set{V}^{+}_v|, K)}{K}, \\
  \bar{\iota}_v
  &= \frac{1}{\binom{|\set{V}^{+}_v|}{2}}
     \sum_{\{i,j\}\subseteq \set{V}^{+}_v}
     \operatorname{IoU}(b_{v,i}, b_{v,j}), \\
  \delta_v
  &= 1 - \bar{\iota}_v,
\end{align}
where \(K\) is the count cap and \(q^{\star}_v=0\) if no candidate has a
non-\text{No} answer. For \(|\set{V}^{+}_v|<2\), set \(\bar{\iota}_v=0\) and
\(\delta_v=1\). Then
\begin{equation}
  r_v =
  \begin{cases}
    \displaystyle
    \frac{
      w_q q^{\star}_v
      + w_s \bar{s}_v
      + w_\kappa \kappa_v
      + w_\delta \delta_v
    }{
      w_q + w_s + w_\kappa + w_\delta
    }, & \set{V}^{+}_v \ne \varnothing, \\[8pt]
    0, & \set{V}^{+}_v = \varnothing.
  \end{cases}
  \label{eq:region-reward}
\end{equation}
The weights are nonnegative and their sum is positive. With the default
implementation they sum to one, but the normalization in
\eqref{eq:region-reward} keeps the reward in \([0,1]\) even if the weights are
changed.

The reward is a heuristic, not a calibrated probability. In particular,
\(q^{\star}_v\) lets one strong candidate dominate, and the diversity term is
maximally favourable when only one box is verified.

% -----------------------------------------------------------------------------
\subsection{Branch selection}

Let \(v\) be an expanded node and let \(\set{C}(v)\) be its active children. If
some child has never been visited, the selector restricts the pool to those
unvisited children. This guarantees that a new region is inspected before the
search commits further to an already promising branch.

If all active children have been visited, the selector computes the cost-aware
UCT index. For each child \(u\in\set{C}(v)\),
\begin{align}
  Q_u &= \frac{R_u}{n_u}, \\
  U_u &= c \sqrt{\frac{\ln(n_v+1)}{n_u}}, \\
  C_u &= \text{remaining-cost forecast below } u, \\
  \mathrm{index}(u)
  &= Q_u + U_u - \lambda \frac{C_u}{T_{\mathrm{ref}}},
\end{align}
and chooses
\begin{equation}
  u^{\star}
  = \argmax_{u \in \set{C}(v)} \mathrm{index}(u).
  \label{eq:selection}
\end{equation}
Ties are broken by the lowest node ID. The three terms have distinct roles:
\begin{itemize}
  \item \(Q_u\) exploits branches that returned high rewards;
  \item \(U_u\) explores branches with few visits;
  \item the last term penalizes branches whose remaining work is expensive.
\end{itemize}
The exploration coefficient is \(c\), the cost weight is \(\lambda\), and
\(T_{\mathrm{ref}}\) normalizes the cost in seconds. Setting \(\lambda=0\)
recovers plain UCT; removing \(U_u\) gives greedy selection.

The remaining-cost forecast is recursive. For a frontier node \(u\),
\begin{equation}
  C_u =
  \begin{cases}
    \max\left(0, \hat{c}(u)\right), & \set{C}(u)=\varnothing, \\[4pt]
    \displaystyle
    \sum_{w\in\set{C}(u)}
      \frac{n_w + 1/|\set{C}(u)|}
           {\sum_{w'\in\set{C}(u)} n_{w'} + 1}
      C_w, & \set{C}(u)\ne\varnothing,
  \end{cases}
  \label{eq:remaining-cost}
\end{equation}
where \(\hat{c}(u)\) estimates the uncached detector, verifier, and split work
at the frontier. The pseudo-visit term \(1/|\set{C}(u)|\) prevents a division by
zero and makes an unvisited set behave like a simple average.

A selected path is built recursively:
\begin{equation}
  P_t = (v_0, v_1, \ldots, v_t),
  \qquad
  v_{\ell+1} = \text{selected child of } v_{\ell},
\end{equation}
until a frontier leaf \(v_t\) is reached. The actual evaluation is performed on
\(v_t\), but the reward is backed up through every node in \(P_t\).

% -----------------------------------------------------------------------------
\subsection{Backup and expansion}

Let \(r_{v_t}\) be the region reward returned by the complete evaluation of the
selected leaf. Every node on the selected path is updated:
\begin{equation}
  n_u \leftarrow n_u + 1,
  \qquad
  R_u \leftarrow R_u + r_{v_t},
  \qquad
  \forall u \in P_t.
  \label{eq:backup}
\end{equation}
The same raw reward is used at every level; there is no discount and no cost
subtraction in the backed-up value. Cost only affects future selection.
Equivalently, the new mean reward is
\begin{equation}
  Q_u \leftarrow
  \frac{n_u Q_u + r_{v_t}}{n_u + 1},
  \qquad
  \forall u \in P_t.
\end{equation}
This is why a promising region increases the selection probability of its entire
branch, not only of the leaf itself.

After the backup, the evaluated node is expanded. Semantic splitting uses a
DINOv3 feature map followed by PCA, SLIC, and K-means:
\begin{equation}
  F_v = \Pi_3\!\left(\phi_{\mathrm{DINO}}(I_v)\right),
\end{equation}
where \(\Pi_3\) projects the patch features to three PCA components. For each
requested SLIC count \(n\in\set{N}\), let \(\ell^{(n)}_v\) be the superpixel
labels and let
\begin{equation}
  H(s) = - \sum_{c} p_{s,c} \log_2 p_{s,c}
\end{equation}
be the colour entropy of superpixel \(s\). The selected count is
\begin{equation}
  n^{\star}
  = \argmin_{n\in\set{N}}
    \frac{1}{|\ell^{(n)}_v|}
    \sum_{s\in \ell^{(n)}_v} H(s).
  \label{eq:slic-entropy}
\end{equation}
K-means then groups superpixels using their mean PCA colour and normalized
centroid:
\begin{equation}
  \mu_s = \left[ \bar{F}_{v,s}, \alpha \bar{p}_s \right],
  \qquad
  k = \min(B, |\ell^{(n^{\star})}_v|),
\end{equation}
where \(B\) is the branching factor and \(\alpha\) is the spatial weight.
Children are the clusters intersected with the parent mask, tightened to their
bounding boxes.

If the resulting children are too few, do not provide a useful zoom, or have
poor geometry, the split falls back to four spatial quadrants. If the node is at
the maximum depth, has too few active pixels, or is too small, it is marked
exhausted and is not split again.

% -----------------------------------------------------------------------------
\subsection{Stopping criteria}

The search stops when any of the following holds:
\begin{equation}
  \begin{aligned}
    &\text{exhausted:} && \text{the root has no unevaluated descendants,} \\
    &\text{api-budget:} && n_{\mathrm{api}} \ge B_{\mathrm{api}}, \\
    &\text{local-budget:} && n_{\mathrm{local}} \ge B_{\mathrm{local}}, \\
    &\text{iterations:} && t \ge T_{\max}, \\
    &\text{match:} && \text{a verified proposal exceeds the match threshold,} \\
    &\text{error:} && \text{a non-recoverable execution error occurred.}
  \end{aligned}
  \label{eq:stop}
\end{equation}
A partial evaluation interrupted by a budget is recorded but is not backed up:
\begin{equation}
  \text{incomplete}(t)
  \quad\Longrightarrow\quad
  \Delta n_u = 0,\ \Delta R_u = 0
  \quad \forall u\in P_t.
\end{equation}
This prevents a failed or truncated evaluation from being treated as a
zero-reward observation.

% -----------------------------------------------------------------------------
\subsection{High-level pseudocode}

\begin{algorithm}[t]
\caption{Bounded cost-aware visual search}
\label{alg:visual-search}
\begin{algorithmic}[1]
\Require image \(I\), question \(q\), configuration \(\theta\)
\Ensure selected region \(R\), candidate boxes \(\set{B}\), tree \(\set{T}\),
         trace \(\tau\)
\State \(s \gets \proc{Prepare}(I,q,\theta)\)
\State \(\set{T} \gets \proc{NewTree}(\proc{Root}(I))\)
\State \(\proc{Expand}(\set{T}.\text{root}, s)\)
\While{not \(\proc{Stop}(\set{T},s)\)}
  \State \(P \gets \proc{SelectPath}(\set{T},s)\)
  \Comment{Eq.~\eqref{eq:selection}}
  \State \(v \gets \text{last}(P)\)
  \State \(r_v \gets \proc{Evaluate}(v,s)\)
  \Comment{Eqs.~\eqref{eq:verified}--\eqref{eq:region-reward}}
  \State \(\proc{Backup}(P,r_v)\)
  \Comment{Eq.~\eqref{eq:backup}}
  \State \(\proc{Expand}(v,s)\)
  \Comment{DINOv3 \(\to\) SLIC \(\to\) K-means, or quadrants}
\EndWhile
\State \(R,\set{B} \gets \proc{Rank}(\set{T})\)
\State \Return \(R,\set{B},\set{T},\tau\)
\end{algorithmic}
\end{algorithm}

The pseudocode keeps the model internals inside \textsc{Evaluate}. The main loop
contains only the four operations that define the search: selection, evaluation,
backup, and expansion. This separation makes the search logic independent of
the particular detector or VLM.

% -----------------------------------------------------------------------------
\subsection{Operations and cost}

Let \(N=T_{\max}\) be the iteration limit, \(D\) the maximum depth, and \(L\) the
maximum number of candidates verified per region. The dominant costs are:
\begin{itemize}
  \item one detector pass per evaluated region, at most \(N\);
  \item at most \(N\) VLM calls in whole-region mode, and at most \(N L\) in
        per-candidate mode, plus repair attempts;
  \item one feature encoding, one SLIC sweep over \(|\set{N}|\) candidate counts,
        and one K-means pass per split;
  \item \(O(|\set{V}|)\) tree memory, with \(|\set{V}|=O(NB)\) in the worst case
        for branching factor \(B\) and depth \(D\).
\end{itemize}
Caches remove repeated work: a node's detector output and split are computed at
most once, and identical verifier inputs can reuse a cached assessment. The tree
can revisit a branch many times, but a node's own complete evaluation runs at
most once; later visits descend into its children.

% -----------------------------------------------------------------------------
\subsection{Summary}

The algorithm can be read as a sequence of simple statements:
\begin{enumerate}
  \item A promising region raises \(Q\) along its path, so the branch is more
        likely to be explored again.
  \item Unvisited children are always considered first, and the UCT bonus
        protects under-visited branches.
  \item The cost penalty removes expensive branches from the top of the queue.
  \item A complete evaluation is the only event that backs up a reward; partial
        or invalid evaluations do not create artificial zeros.
  \item The search returns ranked regions and boxes. The final answer is produced
        downstream by the VLM from the selected region.
\end{enumerate}
 after loading the packages above.
% =============================================================================

\providecommand{\indicator}[1]{\mathbf{1}\!\left[#1\right]}
\providecommand{\set}[1]{\mathcal{#1}}
\providecommand{\argmax}{\operatorname*{arg\,max}}
\providecommand{\argmin}{\operatorname*{arg\,min}}
\providecommand{\proc}[1]{\text{\textsc{#1}}}

\section{Algorithm Walkthrough}
\label{sec:algorithm-walkthrough}

\subsection{High-level idea}

The algorithm is a bounded, best-first tree search over image regions. It does
not build the complete tree in advance. Instead, it starts from the full image
and repeatedly selects one frontier region, evaluates it, and, if the region is
promising and has not reached a stopping rule, splits it into smaller child
regions.

Each iteration therefore has four steps:
\begin{enumerate}
  \item \textbf{Select.} Choose one frontier node using a cost-aware UCT rule.
  \item \textbf{Evaluate.} Run the detector on the selected crop, verify the
        candidate boxes with the VLM, and compute a region reward.
  \item \textbf{Back up.} Add the reward to every node on the selected path.
  \item \textbf{Expand.} Split the evaluated node into child regions, or mark it
        exhausted when no useful split exists.
\end{enumerate}

The intuition is simple. A region that receives a high reward raises the
average reward of its whole branch, so that branch is more likely to be selected
again. At the same time, unvisited children are always explored first, and a UCT
exploration bonus favours branches with few visits. A cost penalty discourages
branches whose remaining work is expensive. The search is therefore an
exploitation--exploration trade-off, not a hard threshold on the VLM score.

The search outputs a selected region \(R\), a set of candidate boxes
\(\set{B}\), and the full tree and allocation trace. In the complete system, the
final free-form answer is generated by the VLM from the selected region:
\begin{equation}
  a = \proc{Answer}(I_R, q),
\end{equation}
where \(I_R\) is the image crop induced by \(R\). The answer is a downstream
product; the tree reward is only a localization heuristic.

% -----------------------------------------------------------------------------
\subsection{Search tree and notation}

Let
\begin{equation}
  \set{T} = (\set{V}, \set{E})
\end{equation}
be the search tree. The root \(v_0\) is the full image region
\begin{equation}
  B_{v_0} = (0, 0, W, H).
\end{equation}
Each node \(v \in \set{V}\) represents a region \(B_v\), and stores
\begin{align}
  n_v &\in \mathbb{N}_0
    && \text{completed evaluations backed up through } v, \\
  R_v &\in \mathbb{R}_{\ge 0}
    && \text{sum of raw rewards backed up through } v, \\
  Q_v &=
    \begin{cases}
      R_v / n_v, & n_v > 0, \\
      0, & n_v = 0,
    \end{cases}
    && \text{empirical mean reward}, \\
  \set{C}(v) &\subseteq \set{V}
    && \text{active children of } v.
\end{align}
A node is \emph{exhausted} when it has no unevaluated descendant within the
depth limit. The frontier is the set of active nodes at which a selected path
can end.

The tree is generated lazily. Initially, only the root exists; it is expanded
once to produce its first children. Every other child is created only after its
parent has been completely evaluated and split. This is important: the algorithm
does not enumerate all possible regions first and inspect them afterwards.

% -----------------------------------------------------------------------------
\subsection{Region evaluation}

We describe the candidate-first path, which is the default. The legacy path
scores the whole region before detection. Let \(v\) be the frontier node
selected at iteration \(t\). The detector is run on the crop
\(I_v = I[B_v]\). It returns a set of candidate detections
\begin{equation}
  \set{D}_v
  = \left\{ (b_{v,i}, s^{\mathrm{det}}_{v,i}) \right\}_{i=1}^{m_v},
\end{equation}
where \(b_{v,i}=(x_{1,i},y_{1,i},x_{2,i},y_{2,i})\) is a box in the crop's
coordinate system and \(s^{\mathrm{det}}_{v,i}\in[0,1]\) is the detector
confidence. Low-confidence boxes are removed, and confidence-ordered non-maximum
suppression keeps a set of distinct candidates
\begin{equation}
  \set{D}^{+}_v
  = \left\{ (b_{v,i}, s^{\mathrm{det}}_{v,i}) \right\}_{i=1}^{m^{+}_v}.
\end{equation}
If \(\set{D}^{+}_v=\varnothing\), the VLM is not called:
\begin{equation}
  r_v = 0.
  \label{eq:empty-region-reward}
\end{equation}
This gate is important for efficiency, because most background regions contain
no candidate at all.

Otherwise, the VLM verifies the candidates, by default against the complete
question \(q\). For candidate \(i\), let
\begin{equation}
  y_{v,i} \in \{\text{Yes}, \text{No}, \text{Uncertain}\},
  \qquad
  \hat{s}^{\mathrm{vlm}}_{v,i} \in [0,1]
\end{equation}
be the verifier answer and score. A candidate is \emph{verified} when
\begin{equation}
  z_{v,i}
  = \indicator{y_{v,i}=\text{Yes}}
    \indicator{\hat{s}^{\mathrm{vlm}}_{v,i} \ge \tau},
  \label{eq:verified}
\end{equation}
where \(\tau\) is the acceptance threshold. Define the checked set and the
verified set as
\begin{equation}
  \set{V}_v = \{1,\ldots,m^{+}_v\},
  \qquad
  \set{V}^{+}_v = \{ i \in \set{V}_v : z_{v,i}=1 \}.
\end{equation}

Candidate ranking uses a fused score
\begin{equation}
  \tilde{s}_{v,i}
  = \frac{
      w_{\mathrm{vlm}} \hat{s}^{\mathrm{vlm}}_{v,i}
      + w_{\mathrm{det}} s^{\mathrm{det}}_{v,i}
    }{
      w_{\mathrm{vlm}} + w_{\mathrm{det}}
    },
  \label{eq:candidate-fusion}
\end{equation}
with \(\hat{s}^{\mathrm{vlm}}_{v,i}=0\) when \(y_{v,i}=\text{No}\). The tree
reward is instead a region-level score, defined next.

% -----------------------------------------------------------------------------
\subsection{Region reward}

The region reward combines the strongest verifier judgment, detector confidence,
the number of verified boxes, and their spatial diversity:
\begin{align}
  q^{\star}_v
  &= \max_{i \in \set{V}_v}
     \left\{
       \hat{s}^{\mathrm{vlm}}_{v,i}
       : y_{v,i} \ne \text{No}
     \right\}, \\
  \bar{s}_v
  &= \frac{1}{|\set{V}^{+}_v|}
     \sum_{i \in \set{V}^{+}_v} s^{\mathrm{det}}_{v,i}, \\
  \kappa_v
  &= \frac{\min(|\set{V}^{+}_v|, K)}{K}, \\
  \bar{\iota}_v
  &= \frac{1}{\binom{|\set{V}^{+}_v|}{2}}
     \sum_{\{i,j\}\subseteq \set{V}^{+}_v}
     \operatorname{IoU}(b_{v,i}, b_{v,j}), \\
  \delta_v
  &= 1 - \bar{\iota}_v,
\end{align}
where \(K\) is the count cap and \(q^{\star}_v=0\) if no candidate has a
non-\text{No} answer. For \(|\set{V}^{+}_v|<2\), set \(\bar{\iota}_v=0\) and
\(\delta_v=1\). Then
\begin{equation}
  r_v =
  \begin{cases}
    \displaystyle
    \frac{
      w_q q^{\star}_v
      + w_s \bar{s}_v
      + w_\kappa \kappa_v
      + w_\delta \delta_v
    }{
      w_q + w_s + w_\kappa + w_\delta
    }, & \set{V}^{+}_v \ne \varnothing, \\[8pt]
    0, & \set{V}^{+}_v = \varnothing.
  \end{cases}
  \label{eq:region-reward}
\end{equation}
The weights are nonnegative and their sum is positive. With the default
implementation they sum to one, but the normalization in
\eqref{eq:region-reward} keeps the reward in \([0,1]\) even if the weights are
changed.

The reward is a heuristic, not a calibrated probability. In particular,
\(q^{\star}_v\) lets one strong candidate dominate, and the diversity term is
maximally favourable when only one box is verified.

% -----------------------------------------------------------------------------
\subsection{Branch selection}

Let \(v\) be an expanded node and let \(\set{C}(v)\) be its active children. If
some child has never been visited, the selector restricts the pool to those
unvisited children. This guarantees that a new region is inspected before the
search commits further to an already promising branch.

If all active children have been visited, the selector computes the cost-aware
UCT index. For each child \(u\in\set{C}(v)\),
\begin{align}
  Q_u &= \frac{R_u}{n_u}, \\
  U_u &= c \sqrt{\frac{\ln(n_v+1)}{n_u}}, \\
  C_u &= \text{remaining-cost forecast below } u, \\
  \mathrm{index}(u)
  &= Q_u + U_u - \lambda \frac{C_u}{T_{\mathrm{ref}}},
\end{align}
and chooses
\begin{equation}
  u^{\star}
  = \argmax_{u \in \set{C}(v)} \mathrm{index}(u).
  \label{eq:selection}
\end{equation}
Ties are broken by the lowest node ID. The three terms have distinct roles:
\begin{itemize}
  \item \(Q_u\) exploits branches that returned high rewards;
  \item \(U_u\) explores branches with few visits;
  \item the last term penalizes branches whose remaining work is expensive.
\end{itemize}
The exploration coefficient is \(c\), the cost weight is \(\lambda\), and
\(T_{\mathrm{ref}}\) normalizes the cost in seconds. Setting \(\lambda=0\)
recovers plain UCT; removing \(U_u\) gives greedy selection.

The remaining-cost forecast is recursive. For a frontier node \(u\),
\begin{equation}
  C_u =
  \begin{cases}
    \max\left(0, \hat{c}(u)\right), & \set{C}(u)=\varnothing, \\[4pt]
    \displaystyle
    \sum_{w\in\set{C}(u)}
      \frac{n_w + 1/|\set{C}(u)|}
           {\sum_{w'\in\set{C}(u)} n_{w'} + 1}
      C_w, & \set{C}(u)\ne\varnothing,
  \end{cases}
  \label{eq:remaining-cost}
\end{equation}
where \(\hat{c}(u)\) estimates the uncached detector, verifier, and split work
at the frontier. The pseudo-visit term \(1/|\set{C}(u)|\) prevents a division by
zero and makes an unvisited set behave like a simple average.

A selected path is built recursively:
\begin{equation}
  P_t = (v_0, v_1, \ldots, v_t),
  \qquad
  v_{\ell+1} = \text{selected child of } v_{\ell},
\end{equation}
until a frontier leaf \(v_t\) is reached. The actual evaluation is performed on
\(v_t\), but the reward is backed up through every node in \(P_t\).

% -----------------------------------------------------------------------------
\subsection{Backup and expansion}

Let \(r_{v_t}\) be the region reward returned by the complete evaluation of the
selected leaf. Every node on the selected path is updated:
\begin{equation}
  n_u \leftarrow n_u + 1,
  \qquad
  R_u \leftarrow R_u + r_{v_t},
  \qquad
  \forall u \in P_t.
  \label{eq:backup}
\end{equation}
The same raw reward is used at every level; there is no discount and no cost
subtraction in the backed-up value. Cost only affects future selection.
Equivalently, the new mean reward is
\begin{equation}
  Q_u \leftarrow
  \frac{n_u Q_u + r_{v_t}}{n_u + 1},
  \qquad
  \forall u \in P_t.
\end{equation}
This is why a promising region increases the selection probability of its entire
branch, not only of the leaf itself.

After the backup, the evaluated node is expanded. Semantic splitting uses a
DINOv3 feature map followed by PCA, SLIC, and K-means:
\begin{equation}
  F_v = \Pi_3\!\left(\phi_{\mathrm{DINO}}(I_v)\right),
\end{equation}
where \(\Pi_3\) projects the patch features to three PCA components. For each
requested SLIC count \(n\in\set{N}\), let \(\ell^{(n)}_v\) be the superpixel
labels and let
\begin{equation}
  H(s) = - \sum_{c} p_{s,c} \log_2 p_{s,c}
\end{equation}
be the colour entropy of superpixel \(s\). The selected count is
\begin{equation}
  n^{\star}
  = \argmin_{n\in\set{N}}
    \frac{1}{|\ell^{(n)}_v|}
    \sum_{s\in \ell^{(n)}_v} H(s).
  \label{eq:slic-entropy}
\end{equation}
K-means then groups superpixels using their mean PCA colour and normalized
centroid:
\begin{equation}
  \mu_s = \left[ \bar{F}_{v,s}, \alpha \bar{p}_s \right],
  \qquad
  k = \min(B, |\ell^{(n^{\star})}_v|),
\end{equation}
where \(B\) is the branching factor and \(\alpha\) is the spatial weight.
Children are the clusters intersected with the parent mask, tightened to their
bounding boxes.

If the resulting children are too few, do not provide a useful zoom, or have
poor geometry, the split falls back to four spatial quadrants. If the node is at
the maximum depth, has too few active pixels, or is too small, it is marked
exhausted and is not split again.

% -----------------------------------------------------------------------------
\subsection{Stopping criteria}

The search stops when any of the following holds:
\begin{equation}
  \begin{aligned}
    &\text{exhausted:} && \text{the root has no unevaluated descendants,} \\
    &\text{api-budget:} && n_{\mathrm{api}} \ge B_{\mathrm{api}}, \\
    &\text{local-budget:} && n_{\mathrm{local}} \ge B_{\mathrm{local}}, \\
    &\text{iterations:} && t \ge T_{\max}, \\
    &\text{match:} && \text{a verified proposal exceeds the match threshold,} \\
    &\text{error:} && \text{a non-recoverable execution error occurred.}
  \end{aligned}
  \label{eq:stop}
\end{equation}
A partial evaluation interrupted by a budget is recorded but is not backed up:
\begin{equation}
  \text{incomplete}(t)
  \quad\Longrightarrow\quad
  \Delta n_u = 0,\ \Delta R_u = 0
  \quad \forall u\in P_t.
\end{equation}
This prevents a failed or truncated evaluation from being treated as a
zero-reward observation.

% -----------------------------------------------------------------------------
\subsection{High-level pseudocode}

\begin{algorithm}[t]
\caption{Bounded cost-aware visual search}
\label{alg:visual-search}
\begin{algorithmic}[1]
\Require image \(I\), question \(q\), configuration \(\theta\)
\Ensure selected region \(R\), candidate boxes \(\set{B}\), tree \(\set{T}\),
         trace \(\tau\)
\State \(s \gets \proc{Prepare}(I,q,\theta)\)
\State \(\set{T} \gets \proc{NewTree}(\proc{Root}(I))\)
\State \(\proc{Expand}(\set{T}.\text{root}, s)\)
\While{not \(\proc{Stop}(\set{T},s)\)}
  \State \(P \gets \proc{SelectPath}(\set{T},s)\)
  \Comment{Eq.~\eqref{eq:selection}}
  \State \(v \gets \text{last}(P)\)
  \State \(r_v \gets \proc{Evaluate}(v,s)\)
  \Comment{Eqs.~\eqref{eq:verified}--\eqref{eq:region-reward}}
  \State \(\proc{Backup}(P,r_v)\)
  \Comment{Eq.~\eqref{eq:backup}}
  \State \(\proc{Expand}(v,s)\)
  \Comment{DINOv3 \(\to\) SLIC \(\to\) K-means, or quadrants}
\EndWhile
\State \(R,\set{B} \gets \proc{Rank}(\set{T})\)
\State \Return \(R,\set{B},\set{T},\tau\)
\end{algorithmic}
\end{algorithm}

The pseudocode keeps the model internals inside \textsc{Evaluate}. The main loop
contains only the four operations that define the search: selection, evaluation,
backup, and expansion. This separation makes the search logic independent of
the particular detector or VLM.

% -----------------------------------------------------------------------------
\subsection{Operations and cost}

Let \(N=T_{\max}\) be the iteration limit, \(D\) the maximum depth, and \(L\) the
maximum number of candidates verified per region. The dominant costs are:
\begin{itemize}
  \item one detector pass per evaluated region, at most \(N\);
  \item at most \(N\) VLM calls in whole-region mode, and at most \(N L\) in
        per-candidate mode, plus repair attempts;
  \item one feature encoding, one SLIC sweep over \(|\set{N}|\) candidate counts,
        and one K-means pass per split;
  \item \(O(|\set{V}|)\) tree memory, with \(|\set{V}|=O(NB)\) in the worst case
        for branching factor \(B\) and depth \(D\).
\end{itemize}
Caches remove repeated work: a node's detector output and split are computed at
most once, and identical verifier inputs can reuse a cached assessment. The tree
can revisit a branch many times, but a node's own complete evaluation runs at
most once; later visits descend into its children.

% -----------------------------------------------------------------------------
\subsection{Summary}

The algorithm can be read as a sequence of simple statements:
\begin{enumerate}
  \item A promising region raises \(Q\) along its path, so the branch is more
        likely to be explored again.
  \item Unvisited children are always considered first, and the UCT bonus
        protects under-visited branches.
  \item The cost penalty removes expensive branches from the top of the queue.
  \item A complete evaluation is the only event that backs up a reward; partial
        or invalid evaluations do not create artificial zeros.
  \item The search returns ranked regions and boxes. The final answer is produced
        downstream by the VLM from the selected region.
\end{enumerate}
 after loading the packages above.
% =============================================================================

\providecommand{\indicator}[1]{\mathbf{1}\!\left[#1\right]}
\providecommand{\set}[1]{\mathcal{#1}}
\providecommand{\argmax}{\operatorname*{arg\,max}}
\providecommand{\argmin}{\operatorname*{arg\,min}}
\providecommand{\proc}[1]{\text{\textsc{#1}}}

\section{Algorithm Walkthrough}
\label{sec:algorithm-walkthrough}

\subsection{High-level idea}

The algorithm is a bounded, best-first tree search over image regions. It does
not build the complete tree in advance. Instead, it starts from the full image
and repeatedly selects one frontier region, evaluates it, and, if the region is
promising and has not reached a stopping rule, splits it into smaller child
regions.

Each iteration therefore has four steps:
\begin{enumerate}
  \item \textbf{Select.} Choose one frontier node using a cost-aware UCT rule.
  \item \textbf{Evaluate.} Run the detector on the selected crop, verify the
        candidate boxes with the VLM, and compute a region reward.
  \item \textbf{Back up.} Add the reward to every node on the selected path.
  \item \textbf{Expand.} Split the evaluated node into child regions, or mark it
        exhausted when no useful split exists.
\end{enumerate}

The intuition is simple. A region that receives a high reward raises the
average reward of its whole branch, so that branch is more likely to be selected
again. At the same time, unvisited children are always explored first, and a UCT
exploration bonus favours branches with few visits. A cost penalty discourages
branches whose remaining work is expensive. The search is therefore an
exploitation--exploration trade-off, not a hard threshold on the VLM score.

The search outputs a selected region \(R\), a set of candidate boxes
\(\set{B}\), and the full tree and allocation trace. In the complete system, the
final free-form answer is generated by the VLM from the selected region:
\begin{equation}
  a = \proc{Answer}(I_R, q),
\end{equation}
where \(I_R\) is the image crop induced by \(R\). The answer is a downstream
product; the tree reward is only a localization heuristic.

% -----------------------------------------------------------------------------
\subsection{Search tree and notation}

Let
\begin{equation}
  \set{T} = (\set{V}, \set{E})
\end{equation}
be the search tree. The root \(v_0\) is the full image region
\begin{equation}
  B_{v_0} = (0, 0, W, H).
\end{equation}
Each node \(v \in \set{V}\) represents a region \(B_v\), and stores
\begin{align}
  n_v &\in \mathbb{N}_0
    && \text{completed evaluations backed up through } v, \\
  R_v &\in \mathbb{R}_{\ge 0}
    && \text{sum of raw rewards backed up through } v, \\
  Q_v &=
    \begin{cases}
      R_v / n_v, & n_v > 0, \\
      0, & n_v = 0,
    \end{cases}
    && \text{empirical mean reward}, \\
  \set{C}(v) &\subseteq \set{V}
    && \text{active children of } v.
\end{align}
A node is \emph{exhausted} when it has no unevaluated descendant within the
depth limit. The frontier is the set of active nodes at which a selected path
can end.

The tree is generated lazily. Initially, only the root exists; it is expanded
once to produce its first children. Every other child is created only after its
parent has been completely evaluated and split. This is important: the algorithm
does not enumerate all possible regions first and inspect them afterwards.

% -----------------------------------------------------------------------------
\subsection{Region evaluation}

We describe the candidate-first path, which is the default. The legacy path
scores the whole region before detection. Let \(v\) be the frontier node
selected at iteration \(t\). The detector is run on the crop
\(I_v = I[B_v]\). It returns a set of candidate detections
\begin{equation}
  \set{D}_v
  = \left\{ (b_{v,i}, s^{\mathrm{det}}_{v,i}) \right\}_{i=1}^{m_v},
\end{equation}
where \(b_{v,i}=(x_{1,i},y_{1,i},x_{2,i},y_{2,i})\) is a box in the crop's
coordinate system and \(s^{\mathrm{det}}_{v,i}\in[0,1]\) is the detector
confidence. Low-confidence boxes are removed, and confidence-ordered non-maximum
suppression keeps a set of distinct candidates
\begin{equation}
  \set{D}^{+}_v
  = \left\{ (b_{v,i}, s^{\mathrm{det}}_{v,i}) \right\}_{i=1}^{m^{+}_v}.
\end{equation}
If \(\set{D}^{+}_v=\varnothing\), the VLM is not called:
\begin{equation}
  r_v = 0.
  \label{eq:empty-region-reward}
\end{equation}
This gate is important for efficiency, because most background regions contain
no candidate at all.

Otherwise, the VLM verifies the candidates, by default against the complete
question \(q\). For candidate \(i\), let
\begin{equation}
  y_{v,i} \in \{\text{Yes}, \text{No}, \text{Uncertain}\},
  \qquad
  \hat{s}^{\mathrm{vlm}}_{v,i} \in [0,1]
\end{equation}
be the verifier answer and score. A candidate is \emph{verified} when
\begin{equation}
  z_{v,i}
  = \indicator{y_{v,i}=\text{Yes}}
    \indicator{\hat{s}^{\mathrm{vlm}}_{v,i} \ge \tau},
  \label{eq:verified}
\end{equation}
where \(\tau\) is the acceptance threshold. Define the checked set and the
verified set as
\begin{equation}
  \set{V}_v = \{1,\ldots,m^{+}_v\},
  \qquad
  \set{V}^{+}_v = \{ i \in \set{V}_v : z_{v,i}=1 \}.
\end{equation}

Candidate ranking uses a fused score
\begin{equation}
  \tilde{s}_{v,i}
  = \frac{
      w_{\mathrm{vlm}} \hat{s}^{\mathrm{vlm}}_{v,i}
      + w_{\mathrm{det}} s^{\mathrm{det}}_{v,i}
    }{
      w_{\mathrm{vlm}} + w_{\mathrm{det}}
    },
  \label{eq:candidate-fusion}
\end{equation}
with \(\hat{s}^{\mathrm{vlm}}_{v,i}=0\) when \(y_{v,i}=\text{No}\). The tree
reward is instead a region-level score, defined next.

% -----------------------------------------------------------------------------
\subsection{Region reward}

The region reward combines the strongest verifier judgment, detector confidence,
the number of verified boxes, and their spatial diversity:
\begin{align}
  q^{\star}_v
  &= \max_{i \in \set{V}_v}
     \left\{
       \hat{s}^{\mathrm{vlm}}_{v,i}
       : y_{v,i} \ne \text{No}
     \right\}, \\
  \bar{s}_v
  &= \frac{1}{|\set{V}^{+}_v|}
     \sum_{i \in \set{V}^{+}_v} s^{\mathrm{det}}_{v,i}, \\
  \kappa_v
  &= \frac{\min(|\set{V}^{+}_v|, K)}{K}, \\
  \bar{\iota}_v
  &= \frac{1}{\binom{|\set{V}^{+}_v|}{2}}
     \sum_{\{i,j\}\subseteq \set{V}^{+}_v}
     \operatorname{IoU}(b_{v,i}, b_{v,j}), \\
  \delta_v
  &= 1 - \bar{\iota}_v,
\end{align}
where \(K\) is the count cap and \(q^{\star}_v=0\) if no candidate has a
non-\text{No} answer. For \(|\set{V}^{+}_v|<2\), set \(\bar{\iota}_v=0\) and
\(\delta_v=1\). Then
\begin{equation}
  r_v =
  \begin{cases}
    \displaystyle
    \frac{
      w_q q^{\star}_v
      + w_s \bar{s}_v
      + w_\kappa \kappa_v
      + w_\delta \delta_v
    }{
      w_q + w_s + w_\kappa + w_\delta
    }, & \set{V}^{+}_v \ne \varnothing, \\[8pt]
    0, & \set{V}^{+}_v = \varnothing.
  \end{cases}
  \label{eq:region-reward}
\end{equation}
The weights are nonnegative and their sum is positive. With the default
implementation they sum to one, but the normalization in
\eqref{eq:region-reward} keeps the reward in \([0,1]\) even if the weights are
changed.

The reward is a heuristic, not a calibrated probability. In particular,
\(q^{\star}_v\) lets one strong candidate dominate, and the diversity term is
maximally favourable when only one box is verified.

% -----------------------------------------------------------------------------
\subsection{Branch selection}

Let \(v\) be an expanded node and let \(\set{C}(v)\) be its active children. If
some child has never been visited, the selector restricts the pool to those
unvisited children. This guarantees that a new region is inspected before the
search commits further to an already promising branch.

If all active children have been visited, the selector computes the cost-aware
UCT index. For each child \(u\in\set{C}(v)\),
\begin{align}
  Q_u &= \frac{R_u}{n_u}, \\
  U_u &= c \sqrt{\frac{\ln(n_v+1)}{n_u}}, \\
  C_u &= \text{remaining-cost forecast below } u, \\
  \mathrm{index}(u)
  &= Q_u + U_u - \lambda \frac{C_u}{T_{\mathrm{ref}}},
\end{align}
and chooses
\begin{equation}
  u^{\star}
  = \argmax_{u \in \set{C}(v)} \mathrm{index}(u).
  \label{eq:selection}
\end{equation}
Ties are broken by the lowest node ID. The three terms have distinct roles:
\begin{itemize}
  \item \(Q_u\) exploits branches that returned high rewards;
  \item \(U_u\) explores branches with few visits;
  \item the last term penalizes branches whose remaining work is expensive.
\end{itemize}
The exploration coefficient is \(c\), the cost weight is \(\lambda\), and
\(T_{\mathrm{ref}}\) normalizes the cost in seconds. Setting \(\lambda=0\)
recovers plain UCT; removing \(U_u\) gives greedy selection.

The remaining-cost forecast is recursive. For a frontier node \(u\),
\begin{equation}
  C_u =
  \begin{cases}
    \max\left(0, \hat{c}(u)\right), & \set{C}(u)=\varnothing, \\[4pt]
    \displaystyle
    \sum_{w\in\set{C}(u)}
      \frac{n_w + 1/|\set{C}(u)|}
           {\sum_{w'\in\set{C}(u)} n_{w'} + 1}
      C_w, & \set{C}(u)\ne\varnothing,
  \end{cases}
  \label{eq:remaining-cost}
\end{equation}
where \(\hat{c}(u)\) estimates the uncached detector, verifier, and split work
at the frontier. The pseudo-visit term \(1/|\set{C}(u)|\) prevents a division by
zero and makes an unvisited set behave like a simple average.

A selected path is built recursively:
\begin{equation}
  P_t = (v_0, v_1, \ldots, v_t),
  \qquad
  v_{\ell+1} = \text{selected child of } v_{\ell},
\end{equation}
until a frontier leaf \(v_t\) is reached. The actual evaluation is performed on
\(v_t\), but the reward is backed up through every node in \(P_t\).

% -----------------------------------------------------------------------------
\subsection{Backup and expansion}

Let \(r_{v_t}\) be the region reward returned by the complete evaluation of the
selected leaf. Every node on the selected path is updated:
\begin{equation}
  n_u \leftarrow n_u + 1,
  \qquad
  R_u \leftarrow R_u + r_{v_t},
  \qquad
  \forall u \in P_t.
  \label{eq:backup}
\end{equation}
The same raw reward is used at every level; there is no discount and no cost
subtraction in the backed-up value. Cost only affects future selection.
Equivalently, the new mean reward is
\begin{equation}
  Q_u \leftarrow
  \frac{n_u Q_u + r_{v_t}}{n_u + 1},
  \qquad
  \forall u \in P_t.
\end{equation}
This is why a promising region increases the selection probability of its entire
branch, not only of the leaf itself.

After the backup, the evaluated node is expanded. Semantic splitting uses a
DINOv3 feature map followed by PCA, SLIC, and K-means:
\begin{equation}
  F_v = \Pi_3\!\left(\phi_{\mathrm{DINO}}(I_v)\right),
\end{equation}
where \(\Pi_3\) projects the patch features to three PCA components. For each
requested SLIC count \(n\in\set{N}\), let \(\ell^{(n)}_v\) be the superpixel
labels and let
\begin{equation}
  H(s) = - \sum_{c} p_{s,c} \log_2 p_{s,c}
\end{equation}
be the colour entropy of superpixel \(s\). The selected count is
\begin{equation}
  n^{\star}
  = \argmin_{n\in\set{N}}
    \frac{1}{|\ell^{(n)}_v|}
    \sum_{s\in \ell^{(n)}_v} H(s).
  \label{eq:slic-entropy}
\end{equation}
K-means then groups superpixels using their mean PCA colour and normalized
centroid:
\begin{equation}
  \mu_s = \left[ \bar{F}_{v,s}, \alpha \bar{p}_s \right],
  \qquad
  k = \min(B, |\ell^{(n^{\star})}_v|),
\end{equation}
where \(B\) is the branching factor and \(\alpha\) is the spatial weight.
Children are the clusters intersected with the parent mask, tightened to their
bounding boxes.

If the resulting children are too few, do not provide a useful zoom, or have
poor geometry, the split falls back to four spatial quadrants. If the node is at
the maximum depth, has too few active pixels, or is too small, it is marked
exhausted and is not split again.

% -----------------------------------------------------------------------------
\subsection{Stopping criteria}

The search stops when any of the following holds:
\begin{equation}
  \begin{aligned}
    &\text{exhausted:} && \text{the root has no unevaluated descendants,} \\
    &\text{api-budget:} && n_{\mathrm{api}} \ge B_{\mathrm{api}}, \\
    &\text{local-budget:} && n_{\mathrm{local}} \ge B_{\mathrm{local}}, \\
    &\text{iterations:} && t \ge T_{\max}, \\
    &\text{match:} && \text{a verified proposal exceeds the match threshold,} \\
    &\text{error:} && \text{a non-recoverable execution error occurred.}
  \end{aligned}
  \label{eq:stop}
\end{equation}
A partial evaluation interrupted by a budget is recorded but is not backed up:
\begin{equation}
  \text{incomplete}(t)
  \quad\Longrightarrow\quad
  \Delta n_u = 0,\ \Delta R_u = 0
  \quad \forall u\in P_t.
\end{equation}
This prevents a failed or truncated evaluation from being treated as a
zero-reward observation.

% -----------------------------------------------------------------------------
\subsection{High-level pseudocode}

\begin{algorithm}[t]
\caption{Bounded cost-aware visual search}
\label{alg:visual-search}
\begin{algorithmic}[1]
\Require image \(I\), question \(q\), configuration \(\theta\)
\Ensure selected region \(R\), candidate boxes \(\set{B}\), tree \(\set{T}\),
         trace \(\tau\)
\State \(s \gets \proc{Prepare}(I,q,\theta)\)
\State \(\set{T} \gets \proc{NewTree}(\proc{Root}(I))\)
\State \(\proc{Expand}(\set{T}.\text{root}, s)\)
\While{not \(\proc{Stop}(\set{T},s)\)}
  \State \(P \gets \proc{SelectPath}(\set{T},s)\)
  \Comment{Eq.~\eqref{eq:selection}}
  \State \(v \gets \text{last}(P)\)
  \State \(r_v \gets \proc{Evaluate}(v,s)\)
  \Comment{Eqs.~\eqref{eq:verified}--\eqref{eq:region-reward}}
  \State \(\proc{Backup}(P,r_v)\)
  \Comment{Eq.~\eqref{eq:backup}}
  \State \(\proc{Expand}(v,s)\)
  \Comment{DINOv3 \(\to\) SLIC \(\to\) K-means, or quadrants}
\EndWhile
\State \(R,\set{B} \gets \proc{Rank}(\set{T})\)
\State \Return \(R,\set{B},\set{T},\tau\)
\end{algorithmic}
\end{algorithm}

The pseudocode keeps the model internals inside \textsc{Evaluate}. The main loop
contains only the four operations that define the search: selection, evaluation,
backup, and expansion. This separation makes the search logic independent of
the particular detector or VLM.

% -----------------------------------------------------------------------------
\subsection{Operations and cost}

Let \(N=T_{\max}\) be the iteration limit, \(D\) the maximum depth, and \(L\) the
maximum number of candidates verified per region. The dominant costs are:
\begin{itemize}
  \item one detector pass per evaluated region, at most \(N\);
  \item at most \(N\) VLM calls in whole-region mode, and at most \(N L\) in
        per-candidate mode, plus repair attempts;
  \item one feature encoding, one SLIC sweep over \(|\set{N}|\) candidate counts,
        and one K-means pass per split;
  \item \(O(|\set{V}|)\) tree memory, with \(|\set{V}|=O(NB)\) in the worst case
        for branching factor \(B\) and depth \(D\).
\end{itemize}
Caches remove repeated work: a node's detector output and split are computed at
most once, and identical verifier inputs can reuse a cached assessment. The tree
can revisit a branch many times, but a node's own complete evaluation runs at
most once; later visits descend into its children.

% -----------------------------------------------------------------------------
\subsection{Summary}

The algorithm can be read as a sequence of simple statements:
\begin{enumerate}
  \item A promising region raises \(Q\) along its path, so the branch is more
        likely to be explored again.
  \item Unvisited children are always considered first, and the UCT bonus
        protects under-visited branches.
  \item The cost penalty removes expensive branches from the top of the queue.
  \item A complete evaluation is the only event that backs up a reward; partial
        or invalid evaluations do not create artificial zeros.
  \item The search returns ranked regions and boxes. The final answer is produced
        downstream by the VLM from the selected region.
\end{enumerate}
